% Préambule commun au rapport complet (rapport.tex) et au rapport final
% (rapport_final.tex). Le fichier qui l'inclut doit définir \ifsuivi avant.

\usepackage[utf8]{inputenc}
\usepackage[T1]{fontenc}
\usepackage[french]{babel}
\usepackage{lmodern}
\usepackage[hmargin=1.8cm, top=1.2cm, bottom=1cm, includehead, includefoot, headsep=8pt, footskip=16pt]{geometry}
\usepackage{amsmath}
\usepackage{graphicx}
\usepackage{booktabs}
\usepackage[section]{placeins} % les tableaux restent dans leur section
\usepackage{tabularx}
\usepackage{array}
\usepackage{enumitem}
\usepackage[dvipsnames]{xcolor}
\usepackage{siunitx}
\usepackage{fancyhdr}
\usepackage[font=small, labelfont=bf]{caption}
\usepackage[hidelinks]{hyperref}

\sisetup{locale = FR}
\graphicspath{{figures/}}
\setlist{nosep, leftmargin=*}
\newcolumntype{Y}{>{\raggedright\arraybackslash}X}
\renewcommand{\arraystretch}{1.15}

% Encadré de travail, visible seulement en version de suivi. Contenu : une liste
% \item ... de ce qui reste à faire dans la section.
\newcommand{\afaire}[1]{\ifsuivi\par\medskip\noindent
  \fcolorbox{BrickRed}{BrickRed!6}{\parbox{\dimexpr\linewidth-2\fboxsep-2\fboxrule}{\small
    \textcolor{BrickRed}{\textbf{À faire}}
    \begin{itemize}#1\end{itemize}}}\par\medskip\fi}
% Renvoi vers le code qui produit les chiffres d'une partie (version de suivi seulement).
\newcommand{\outrouver}[1]{\ifsuivi\par\smallskip\noindent
  {\small\sloppy\color{NavyBlue}\textbf{Où trouver :} #1\par}\smallskip\fi}
% Nom de colonne / de fichier du dépôt.
\newcommand{\code}[1]{\texttt{\detokenize{#1}}}
% Autorise une coupure de ligne après / et - dans les noms longs (sans tiret ajouté).
% Attention : les espaces sont ignorés ; pour un nom avec espace, utiliser \texttt.
\renewcommand{\code}[1]{\texttt{\codecut#1\codeend}}
\def\codeend{\codeend}
\def\codecut#1{\ifx#1\codeend\else\ifx#1/\slash\else\ifx#1-\char`\-\allowbreak\else\detokenize{#1}\fi\fi\expandafter\codecut\fi}
